\ifdefined\iscombined
	\lecturetitle{Relations}
\else
\documentclass{article}
\usepackage{listings}
\usepackage{amsthm}
\usepackage{comment}
\usepackage{amsmath}
\usepackage{amsfonts}
\usepackage{sectsty}
\usepackage{hyperref}
\usepackage{fancyhdr}
\usepackage[top=1in,bottom=1in,left=1in,right=1in]{geometry}
\usepackage{float}
\usepackage{graphicx}
\usepackage{tabularx}
\usepackage{mathtools}

\date{
	\vspace{-.75em}
	{\large \today}
}

\title{
	\vspace*{-1.5em}
	{\Huge Lecture 8} \\
	\vspace{-1em}
	\rule{\textwidth/4}{.01em} \\
	\vspace{-.25em}
	{\Large Relations}
	\vspace{-.75em}
}

\author{
	{\large Matthew Farah}
}

\hypersetup{
	colorlinks=true,
	linkcolor=blue,
	filecolor=magenta,
	urlcolor=blue,
	citecolor=blue,
	anchorcolor=blue
}

\fancyhead{}
\fancyhead[L]{\today}
\fancyhead[R]{Lecture 8 - Relations}

\sectionfont{\fontsize{12}{12}\bfseries}
\subsectionfont{\fontsize{10}{12}\normalfont}
\subsubsectionfont{\fontsize{10}{12}\normalfont}

\begin{document}

\maketitle
\pagestyle{fancy}

\fi

\begin{itemize}
	\item Fault injection involves deliberately introducing faults into the system.
	\begin{itemize}
		\item Done to evaluate efficacy of existing tests.
		\begin{itemize}
			\item Kinda like coverage metrics.
		\end{itemize}
		\item More expensive than coverage metrics.
	\end{itemize}
	\item Fault injection can be done in one of two ways:
	\begin{enumerate}
		\item Compile-time fault injection.
		\begin{itemize}
			\item Introduces faults into the source code.
			\item Typically done manually.
			\item Performed in three possible ways:
			\begin{enumerate}
				\item Operand replacement.
				\begin{itemize}
					\item Ex. Changing $a + b$ to $a + c$.
				\end{itemize}
				\item Expression modification.
				\begin{itemize}
					\item Ex. Changing $a + b$ to $a - b$.
				\end{itemize}
				\item Statement modification.
				\begin{itemize}
					\item Ex. Deleting $x \coloneq 1$ from code.
				\end{itemize}
			\end{enumerate}
		\end{itemize}
		\item Runtime fault injection.
		\begin{itemize}
			\item Introduces faults as the program is running.
			\item Typically automated.
			\item Typically performed by use of special software triggers.
		\end{itemize}
	\end{enumerate}
	\item A binary relation is a set $R$ representing ordered associations between some sets $A$ and $B$.
	\begin{itemize}
		\item $R \subseteq A \times B$.
		\item Elements are expressed as $(a, b) \in R$.
		\begin{itemize}
			\item $a \in A, b \in B$.
		\end{itemize}
	\end{itemize}
	\item Equivalence relations are special kinds of binary relations.
	\item Equivalence relations $R \subseteq A \times A$ must satisfy the following three properties:
	\begin{enumerate}
		\item Reflexivity.
		\begin{itemize}
			\item $\forall a \in A; (a, a) \in R$.
		\end{itemize}
		\item Symmetry.
		\begin{itemize}
			\item $\forall a_1, a_2 \in A; (a_1, a_2) \in R \implies (a_2, a_1) \in R$.
		\end{itemize}
		\item Transitivity.
		\begin{itemize}
			\item $\forall a_1, a_2, a_3 \in A; (a_1, a_2) \in R \wedge (a_2, a_3) \in R \implies (a_1, a_3) \in R$.
		\end{itemize}
	\end{enumerate}
	\item Binary and equivalence relations are fundamental to test design theory.
\end{itemize}

\ifdefined\iscombined
	\newpage
\else
	\end{document}
\fi
